% Part: lambda-calculus
% Chapter: syntax
% Section: de-bruijn

\documentclass[../../../include/open-logic-section]{subfiles}

\begin{document}

\olfileid[pt]{lam}{syn}{deb}

\olsection{O índice de De Bruijn}

A equivalência~$\alpha$ é muito natural, pois termos
$\alpha$-equivalentes «significam a mesma coisa». De fato, é possível dar
uma sintaxe para termos lambda que não distingue termos que podem ser
convertidos um no outro por~$\alpha$. A mais conhecida substitui as variáveis
por seu \emph{índice de De Bruijn}.

Quando escrevemos $\lambd[x][M]$, declaramos explicitamente que $x$ é o
parâmetro da função, de modo que possamos usar $x$ em $M$ para nos referirmos
a esse parâmetro. No índice de De Bruijn, contudo, os parâmetros não têm nome,
e a referência a eles no corpo da função é denotada por um número que indica
os níveis de abstração entre ambos. Por exemplo, considere
$\lambd[x][\lambd[y][y x]]$: a abstração externa liga a variável~$x$; a
abstração interna liga a variável~$y$; o subtermo $y x$ está no escopo da
abstração interna: não há abstração entre $y$ e sua abstração~$\lambd[y]$,
mas há uma abstração entre $x$ e sua abstração~$\lambd[x]$. Assim, escrevemos
$0\, 1$ para $y x$, e $\lambd[][\lambd[][01]]$ para o termo inteiro.

\begin{defn}
  Os termos de De Bruijn são definidos indutivamente da seguinte maneira:
  \begin{enumerate}
  \item $n$, onde $n$ é qualquer número natural.
  \item $PQ$, onde $P$ e $Q$ são ambos termos de De Bruijn.
  \item $\lambd[][N]$, onde $N$ é um termo de De Bruijn.
  \end{enumerate}
\end{defn}

Uma tradução formal de termos lambda ordinários para termos indexados de De
Bruijn é dada a seguir:
\begin{defn}
  \begin{align*}
    F_\Gamma(x) &= \Gamma(x) \\
    F_\Gamma(PQ) &= F_\Gamma(P)F_\Gamma(Q) \\
    F_\Gamma(\lambd[x][N]) &= \lambd[][F_{x,\Gamma}(N)]
  \end{align*}
  onde $\Gamma$ é uma lista de variáveis indexada a partir de zero, e
  $\Gamma(x)$ denota a posição da variável $x$ em $\Gamma$. Por exemplo, se
  $\Gamma$ é $x,y,z$, então $\Gamma(x)$ é $0$ e $\Gamma(z)$ é $2$.
  
  $x,\Gamma$ denota a lista resultante da inserção de $x$ no início de
  $\Gamma$; por exemplo, continuando com o $\Gamma$ do último exemplo,
  $w,\Gamma$ é $w,x,y,z$.
\end{defn}

A recuperação de um termo lambda padrão a partir de um termo de De Bruijn é
feita da seguinte maneira:

\begin{defn}
  \begin{align*}
    G_\Gamma(n) &= \Gamma[n] \\
    G_\Gamma(PQ) &= G_\Gamma(P) G_\Gamma(Q) \\
    G_\Gamma(\lambd[][N]) &= \lambd[x][G_{x,\Gamma}(N)]
  \end{align*}
  onde $\Gamma$ é novamente uma lista de variáveis indexada a partir de zero,
  e $\Gamma[n]$ denota a variável na posição~$n$. Por exemplo, se $\Gamma$ é
  $x,y,z$, então $\Gamma[1]$ é $y$.

  Na última equação, escolhe-se para $x$ qualquer variável que não pertença a
  $\Gamma$.
\end{defn}

Apresentamos aqui alguns resultados sem demonstrá-los:

\begin{prop}
  Se $M \aconvone M'$ e $\Gamma$ é qualquer lista que contenha $\FV{M}$,
  então $F_\Gamma(M) \eqs F_\Gamma(M')$.
\end{prop}

\end{document}
